Le comptage précis et en temps réel de petits éléments de fixation métalliques sur des bandes transporteuses industrielles à haute vitesse demeure un goulot d'étranglement persistant dans la logistique de production. L'inspection visuelle manuelle est éprouvante sur le plan cognitif, sujette aux erreurs et économiquement insoutenable, tandis que les méthodes conventionnelles d'estimation par pesée manquent de la granularité requise pour la traçabilité moderne des chaînes d'approvisionnement. Les solutions automatisées existantes sont en outre entravées par les reflets spéculaires sur les surfaces réfléchissantes, le chevauchement dense des pièces et le coût de calcul prohibitif des détecteurs de haute précision des défis qui excluent collectivement tout déploiement fiable et permanent en périphérie d'usine.\\

Ce mémoire présente un framework edge-AI de bout en bout, indépendant du cloud, pour la détection, la classification et le comptage en temps réel de fixations métalliques sur convoyeurs industriels. Au cœur du système se trouve YOLO11-GED, une architecture légère de détection d'objets enrichie d'un nouveau module EMAc2f qui combine de manière synergique trois mécanismes complémentaires : l'attention multiple efficace (Efficient Multi-Attention) pour la suppression des artefacts spéculaires, la réutilisation de caractéristiques fondée sur GhostNet pour une représentation des petits objets économe en paramètres, et le suréchantillonnage dynamique (Dynamic Upsampling) pour préserver la fidélité spatiale à un coût de calcul réduit. Afin d'établir une base de référence rigoureuse, neuf configurations de détecteurs couvrant trois générations de YOLO (YOLOv8n/s/m, YOLOv11n/s/m, YOLOv26n/s/m) sont évaluées sur un jeu de données spécialement constitué d'images de fixations capturées dans des conditions réalistes de convoyage, incluant un éclairage variable, un flou de mouvement et des amas denses.\\

Le modèle entraîné est déployé sur une plateforme embarquée NVIDIA Jetson avec accélération TensorRT et interfacé avec une caméra de profondeur Intel RealSense, atteignant un objectif de mAP@0,5 ≥ 0,95 à ≥ 30 IPS sans aucune dépendance au cloud. Autour du moteur de vision, une plateforme de production full-stack est mise en œuvre, comprenant un service d'inférence Python, une API REST Express.js, une base de données persistante pour les journaux d'audit, une couche Redis pour la télémétrie en temps réel, ainsi que deux interfaces opérateur (tableau de bord web et application de bureau). Le système intégré est validé sur un banc d'essai physique de convoyeur, démontrant un comptage continu à faible latence qui élimine les recomptages manuels et réduit l'erreur humaine.\\

Au-delà du livrable industriel, ce travail contribue à une méthodologie reproductible de la conception au déploiement, à des études d'ablation quantitatives isolant l'effet de chaque composant d'EMAc2f, et à des preuves documentées que des modifications architecturales ciblées permettent d'améliorer le rappel des petits objets sous des budgets de calcul embarqué stricts, sans sacrifier le débit.